\section{Tres AI frontier: generación, acción y control físico}

Los recursos de base se incorporan a distintas formas de aplicación. La clase divide el enfoque en generative, agentic y physical AI; el valor didáctico de esta clasificación es que el failure mode y la evidence de cada una son completamente distintos. El contenido generado puede evaluarse fuera de línea; un agent modifica el workflow state; un physical system afecta directamente a personas y equipos, por lo que los requisitos de validación, permisos y recuperación aumentan de forma escalonada.

\subsection{De la calidad de la salida a la seguridad de la acción}

La sección anterior explicó que los tres tipos de sistemas de IA tienen distintos radios de acción; esta sección convierte esa diferencia en requisitos de validación. Al leer la figura, observe el orden generative, agentic, physical: el sistema se amplía de «generar candidatos» a «modificar el estado digital», y luego a «controlar el mundo físico»; cada paso aumenta los permisos, el impacto, la demora en la detección y el costo de recuperación, por lo que no puede reutilizarse el mismo benchmark.

\lecturefigure{07-three-ai-frontiers.png}{Generative, agentic y physical AI requieren evidencia y rendición de cuentas progresivamente más estrictas.}{Subtítulos locales 16:10--20:20; redibujo conceptual.}

\begin{knowledgebox}{Lectura de la figura: un mismo benchmark no puede validar los tres niveles de sistema}
En generative AI pueden revisarse factuality, utility y diversity; en agentic AI también hay que revisar tool correctness, state transition, retry y exception handling; en physical AI es indispensable agregar sensor failure, control stability, safe stop, human override y la normativa aplicable en sitio. La puntuación del modelo es solo una parte de la evidencia del sistema.
\end{knowledgebox}

\begin{importantbox}{A mayor radio de acción, mayor costo de validación}
El riesgo puede entenderse de forma aproximada como
\[
R \propto P(failure)\times I(impact)\times D_{detect},
\]
donde $P(failure)$ es la probabilidad de falla, $I(impact)$ es el impacto y $D_{detect}$ es la demora en detectar y corregir. Cuando un Agent obtiene más permisos, aunque la tasa de error del modelo no cambie, el impacto y la demora de recuperación pueden amplificarse; por ello hay que agregar minimización de permisos, simulation, approval gate y rollback.
\end{importantbox}

\subsection{Salud: acelerar la búsqueda no significa que desaparezca la gobernanza}

La clase usa las formulaciones de fármacos y el trasplante de corazón robótico para ilustrar la entrada de la IA en la medicina. Los materiales oficiales del KFSHRC confirman que en 2024 se realizó un fully robotic heart transplant; aquí, «fully robotic» describe la ruta quirúrgica y el uso del robotic system, y no debe reescribirse como una autonomous surgery sin un médico responsable. De manera similar, la IA puede acortar la candidate generation y la formulation search, pero el clinical trial, la regulatory review y la evidencia de seguridad a largo plazo siguen existiendo.

\lecturefigure{08-healthcare-boundaries.png}{La IA médica puede comprimir la búsqueda en investigación, pero no puede omitir la evidencia clínica ni los límites de responsabilidad.}{Subtítulos locales 16:40--18:20; materiales oficiales del KFSHRC sobre el trasplante de corazón totalmente robótico.}

\begin{knowledgebox}{Lectura de la figura: no sumar las dos líneas de tiempo}
A la izquierda está el research loop: generar candidatos, predecir propiedades y programar experimentos; a la derecha, la care delivery: selección de pacientes, planificación quirúrgica, control médico y monitoreo posoperatorio. La IA puede mejorar algunos pasos de ambas cadenas, pero cada cadena tiene datos, ética, regulación y responsables distintos.
\end{knowledgebox}

\begin{warningbox}{«De diez años a dos» debe entenderse como un caso de clase}
Las cifras de duración que aparecen en los subtítulos no cuentan con material suficiente para confirmar su punto de inicio, su punto final ni el proyecto de enfermedad. Las notas conservan solo la conclusión transferible: la IA puede acortar la búsqueda de formulaciones y experimentos, pero eso no permite afirmar que el ciclo completo de drug development o de approval se comprima en la misma proporción.
\end{warningbox}

\subsection{Energía: el Agent debe integrarse al ciclo cerrado, no solo generar recomendaciones}

La industria energética dispone de sensor, maintenance, geology y operational history, y también tiene altos costos de seguridad. El ponente mencionó la corrosion y el drilling; el verdadero problema de sistema es cómo entra la recomendación del modelo en una work order, quién la aprueba, cómo regresan los resultados y cómo se detiene la operación ante una anomalía. Un agent sin outcome capture solo puede generar análisis y no logra un operational learning sostenible.

\lecturefigure{09-energy-agent-loop.png}{Un Agent vertical necesita que datos, recomendaciones, decisiones humanas y retroalimentación de resultados formen un ciclo cerrado.}{Subtítulos locales 18:20--20:20; redibujo conceptual.}

\begin{knowledgebox}{Lectura de la figura: por qué la domain data sigue sin ser suficiente}
Los datos históricos pueden reflejar equipos antiguos, estrategias anteriores y selection bias; la falta de datos de sensores puede producir false confidence; las operation exitosas a menudo carecen de etiquetas detalladas. Por ello, el Agent necesita uncertainty, operator rationale, counterfactual e incident review, en lugar de entregar la proprietary data directamente a un modelo de propósito general.
\end{knowledgebox}

\teachervoice{La clase considera las capacidades industriales existentes de un país como punto de partida de los AI use case: un país energético no necesita entrenar solo modelos «grandes y generalistas», sino que puede construir primero ventajas en workflow de alto valor, alta densidad de datos y con retroalimentación. Esto se acerca más a la ingeniería de sistemas que «primero tener el modelo y luego buscar el problema».}

\subsection{Resumen de la sección}

Las tres AI frontier comparten una base común de modelos y capacidad de cómputo, pero la diferencia real está en el radio de acción, el ciclo cerrado de datos y la estructura de responsabilidades. Cuanto más cerca se está de los servicios públicos y del mundo físico, menos puede un único benchmark representar la madurez de un sistema.
